% id: 246-584
% book: 개념원리 공통수학2 · 29 역함수
% section: 연습문제 실력 UP
% figure: none
% answer: ①
% answer_source: 답지
% variant_level: 0 (원본 전사)
% 이 파일은 build.mjs 가 items.json 에서 만든다. 고칠 때는 items.json 을 고친다.
\smash{\raisebox{1.5mm}{\dmkichul{개념원리 공통수학2 246쪽 584번 · 교육청 기출}}}
세 집합 \par\smallskip\dispeq{$X=\{1{,}\mkern3mu\mathopen{}2{,}\mkern3mu\mathopen{}3{,}\mkern3mu\mathopen{}4\}$, $Y=\{2{,}\mkern3mu\mathopen{}3{,}\mkern3mu\mathopen{}4{,}\mkern3mu\mathopen{}5\}$, $Z=\{3{,}\mkern3mu\mathopen{}4{,}\mkern3mu\mathopen{}5\}$}\par\smallskip\noindent 에 대하여 두 함수 $f:X\longrightarrow Y$, $g:Y\longrightarrow Z$가 다음 조건을 만족시킨다. \exprbox{\begin{minipage}{0.96\linewidth}\raggedright ㈎ 함수 $f$는 일대일대응이다.\\ ㈏ $x\in(X\cap Y)$이면 $g(x)-f(x)=1$이다.\end{minipage}}\noindent 보기에서 옳은 것만을 있는 대로 고른 것은? \bogi{ㄱ. 함수 $g\circ f$의 치역은 $Z$이다.\\ ㄴ. $f^{-1}(5)\ge 2$\\ ㄷ. $f(3)<g(2)<f(1)$이면 $f(4)+g(2)=6$이다.}
\begin{choices32}{ㄱ}{ㄱ, ㄴ}{ㄱ, ㄷ}{ㄴ, ㄷ}{ㄱ, ㄴ, ㄷ}\end{choices32}
